\section{Roadmap}
\label{sec:roadmap}

This section reorganizes the backlog in \texttt{docs/TODO.md} (items R1--R7
below map directly to that file's items 1--7) into three phases by
dependency and cost, and adds the two items raised elsewhere in this report
but not present in \texttt{TODO.md}. No calendar dates are attached beyond
``resume in September'' -- phases reflect priority and dependency order, not
a schedule, since real-world constraints on the September restart are not
yet known.

\begin{figure}[H]
\centering
\begin{tikzpicture}[
  every node/.style={font=\small},
  phase/.style={rectangle, rounded corners, draw=rrblue, fill=rrblue!8,
    minimum width=4.3cm, minimum height=2.6cm, align=left, text width=4cm,
    inner sep=8pt},
  phasehead/.style={font=\bfseries\small\color{rrblue}}
]
  \node[phase] (a) at (0,0) {
    \textbf{Phase A -- Verify \& stabilize}\\[2pt]
    \footnotesize
    $\bullet$ Re-verify Session 5 geometry fixes on hardware (I1) -- done\\
    $\bullet$ Re-check camera/board calibration (I8) -- done\\
    $\bullet$ R1: zone-bounds-aware release placement -- done
  };
  \node[phase, right=0.9cm of a] (b) {
    \textbf{Phase B -- Model quality \& schema}\\[2pt]
    \footnotesize
    $\bullet$ R4: bounding-box scene schema -- done\\
    $\bullet$ R5: rewrite skill/primitive prompt text\\
    $\bullet$ R3: real tray-dimension-aware placement\\
    $\bullet$ Benchmark: add VLM\_LLM condition -- attempted, not complete
  };
  \node[phase, right=0.9cm of b] (c) {
    \textbf{Phase C -- New capabilities}\\[2pt]
    \footnotesize
    $\bullet$ R2: push/nudge skill (no grasp)\\
    $\bullet$ R6: multi-call ``chat'' interaction\\
    $\bullet$ R7: geometry fixes as callable skills
  };

  \draw[-{Latex[length=3mm]}, thick, rrblue!50] (a.east) -- (b.west);
  \draw[-{Latex[length=3mm]}, thick, rrblue!50] (b.east) -- (c.west);
\end{tikzpicture}
\caption{Proposed phasing of the roadmap. Phase A is the recommended
  September restart focus; B depends on A's schema decisions; C depends on
  R6/R7 being designed together (Section~\ref{sec:roadmap}, R6).}
\label{fig:roadmap-phases}
\end{figure}

\subsection{Phase A -- verify and stabilize}

\paragraph{Hardware re-verification (I1, I8).} \textbf{Done.} The September
120-trial benchmark (Section~\ref{sec:benchmark}) exercised both the
Session 5 geometry fixes and the camera/calibration path extensively on
live hardware, not a replayed/synthetic check -- 100\% task safety across
all 120 trials is the strongest verification either could have gotten.
Kept here, marked done, rather than removed, since this paragraph was the
gate the rest of this roadmap was written to wait on.

\paragraph{R1 -- Zone-bounds-aware release placement.} \textbf{Done,
September 2026 (ROBOAI-18)}, resolved as a side effect of a hardware bug
found during the September benchmark rather than planned as its own session
-- see Section~\ref{sec:tech-decisions}. What follows is the original
pre-implementation plan, kept for context.
\texttt{distribute\_zone\_releases} (\texttt{schemas.py}) currently spaces
out colliding releases along a line/grid using only the released object's
\texttt{grasp\_width}, with no notion of how big the target tray/zone
actually is -- nothing stops the grid from growing past the zone's physical
edge, and \texttt{PlanValidator} only checks \emph{global} workspace limits,
not per-zone bounds, so an overflow would not be caught anywhere. Zone size
is reliably known in \texttt{LLM} mode (\texttt{scene\_mock.json}'s
hand-authored \texttt{targets.*.size}) and \texttt{VLM\_LLM} mode (the
generated scene's visually-estimated target size), but not in plain
point-grounding \texttt{VLM} mode. Plan: add \emph{opportunistic}
bounds-checking -- clamp/validate against a size when one is available, fall
back to today's unbounded behavior when it is not. Needs re-plumbing an
optional size hint back into \texttt{distribute\_zone\_releases}, which
currently deliberately takes only \texttt{plan} (no \texttt{scene\_state}) --
see Section~\ref{sec:tech-decisions}'s zone-collision-spacing discussion for
why that dependency was dropped, and keep any size hint optional so it does
not reintroduce the LLM-only-matching problem that fix solved.

\subsection{Phase B -- model quality and schema}

\paragraph{R4 -- \texttt{scene\_mock.json}: bounding box instead of
position + size, for \texttt{targets}.} \textbf{Done, September 2026
(ROBOAI-19).} What follows is the original pre-implementation plan, kept
for context; Section~\ref{sec:tech-decisions} has what was actually built,
including one correction the plan below got wrong (the stacking-height
instruction was duplicated in only two of six prompt files, not all six --
the other four had none at all, masked by the old schema letting the model
just copy a ready-made \texttt{position}). Replace
\texttt{targets.*.position} (center) + \texttt{size} (dimensions) with an
explicit bounding box (\texttt{min}/\texttt{max} corners, or a \texttt{min}
corner + size). Agreed and deliberately deferred until after the benchmark
push (now complete). Rationale: removes an assumed, undocumented convention
(\texttt{position} = center, half-extent = \texttt{size}/2) that this
project's own bugs have repeatedly tripped on (the release-height
base-vs-contact-point confusion, the footprint-vs-point-distance matching
gap in \texttt{\_fix\_release\_height}); matches how a physical zone is
actually measured by hand (two corners); simplifies the existing
footprint-containment code (direct min/max comparison instead of computing
half-extents every time) and the stacking-height prompt instruction
(\texttt{target.z\_max} instead of
\texttt{target.position.z + target.size[2]}); and fits VLM bbox grounding
more naturally (a detected pixel bbox's corners deproject straight into a
real-world min/max box). \textbf{Not} proposed for \texttt{objects.*} --
picking needs a contact point, not a box. \textbf{Breaking change}, touches:
\texttt{llm\_planner\_node.py}'s \texttt{\_fix\_object\_height}/
\texttt{\_fix\_release\_height}; the stacking-height instruction duplicated
across all six LLM prompt files; \texttt{vlm\_llm\_planner\_node.py}'s
\texttt{\_build\_generated\_scene}; \texttt{docs/GRASP\_GEOMETRY\_PIPELINE.md};
and \texttt{scene\_mock.json} itself.

\paragraph{R5 -- Rewrite the skill/primitive descriptions for clarity.}
Revisit what the LLM/VLM actually sees: the \texttt{UR5Action} field
docstrings (\texttt{extraction\_classes.py}) and the skill-list text embedded
in each reasoning method's prompt (\texttt{skills.py} / per-method
\texttt{EmbodiedAgentsPrompts/*.py}), to better explain what each primitive
does, what each parameter means, and what convention it follows -- e.g.\ is
\texttt{object\_height} still worth exposing to the model at all, now that
\texttt{\_fix\_object\_height}/\texttt{\_fix\_release\_height} compute it
deterministically and ignore whatever the model puts there (closely related
to R7). No concrete plan yet; needs a pass through current prompt text to
find what is unclear or stale first.

\paragraph{R3 -- Real tray-dimension-aware ``best'' positioning.} Distinct
from R1's bounds-\emph{checking}: actively using a tray's real physical size
to compute the \emph{best} placement, not just a safe non-colliding one --
centering the layout, choosing the axis that best fits the tray's actual
aspect ratio instead of always growing along a fixed axis (currently
hardcoded +y then +x, see \texttt{\_zone\_slot\_offset}), or packing more
optimally than a simple line/grid. Worth deciding whether this stays
deterministic (inside \texttt{distribute\_zone\_releases}, as today's
line/grid does) or becomes an explicit skill the LLM can invoke with a
stated intent (``arrange neatly'' vs.\ ``just don't collide''). Depends on
R1 first, since both need real zone dimensions plumbed through.

\paragraph{Benchmark extension.} \textbf{Attempted September 2026, not
completed.} A third \texttt{VLM\_LLM} arm was added to the September rerun
of the 120-trial study (Section~\ref{sec:benchmark}) but abandoned after 3
trials: a real object in the camera's field of view returned no usable
depth at its viewing angle and crashed the hybrid grounding call every time
it was detected -- a physical-setup problem, not evidence about the hybrid
architecture itself. Revisiting this is now mostly a matter of clearing
that object from the workspace and re-running, not further design work.

\subsection{Phase C -- new capabilities}

\paragraph{R2 -- Non-grasping ``push/nudge'' skill.} A skill that
repositions an object by contacting and sliding it with the end-effector,
without closing the gripper -- e.g.\ nudging into place, clearing a path, or
repositioning something the gripper cannot or should not grasp. Open
questions before implementing: the action shape (likely a new
\texttt{ALLOWED\_SKILLS} entry alongside
\texttt{approach|pick|release|move\_home|wait} in \texttt{schemas.py}, needing
a start-contact point and an end point); the executor-side trajectory
(contact-and-slide -- different geometry from pick's mid-body descent); and
threading the new skill through every reasoning method's prompt (LLM + VLM
variants), the same way \texttt{grasp\_width} and bbox grounding were
threaded through everywhere in earlier passes.

\paragraph{R6 -- A ``chat'' for multiple LLM/VLM calls per task.} Every
reasoning method today is a single request/response cycle (or, for
\texttt{cot\_sc}, several \emph{independent} samples voted on at once) -- no
notion of an ongoing conversation where a follow-up call can react to what
an earlier call in the \emph{same} task decided. Explore a real multi-turn
interaction, closer to ReAct's reason/act loop but generalized, so the agent
(or the operator) can issue several calls within one task and have later
calls see earlier ones' output/state. Directly enables R7 (a model that can
call a ``fix my plan'' skill and see the result before finalizing).

\paragraph{R7 -- Expose the deterministic geometry fixes as callable
skills.} Several deterministic, code-side corrections currently silently
override whatever the LLM/VLM computed:
\texttt{llm\_planner\_node.py}'s \texttt{\_fix\_object\_height},
\texttt{\_fix\_release\_height} (surface-height lookup, footprint-overlap
nudging), and \texttt{schemas.py}'s \texttt{distribute\_zone\_releases}
(collision spacing). All are one-shot, invisible post-processing today: the
model never sees that its numbers were corrected and cannot ask for a
\emph{different} correction if the automatic one is not what was actually
wanted. Idea: expose them as callable tools the LLM/VLM can invoke directly
-- ``recompute this release's height'', ``space these releases apart'',
``nudge this point out of zone X'' -- as part of a multi-call interaction
(R6), so the model actively decides to use the correction rather than having
it silently applied. Would also make failures more diagnosable (the model's
own trace would show \emph{why} a position changed). Open question: how much
of the current hardcoded-in-Python geometry logic to keep as an automatic
safety net regardless -- this project's whole pattern so far has been ``do
not trust the model's own arithmetic for real-world geometry''
(Section~\ref{sec:tech-decisions}) -- versus how much to hand over as an
optional tool call.

\subsection{Not on this roadmap, deliberately}

Two things this report does \emph{not} recommend picking back up without a
separate decision: reviving the VLA direction (Section~\ref{sec:introduction}
explains why it was abandoned; nothing in the benchmark results
(Section~\ref{sec:benchmark}) argues for revisiting that), and adding new
foundation-model providers beyond the current five (Groq, Nebius, OpenAI,
Anthropic, Gemini) -- the benchmark's findings point to reasoning/arithmetic
robustness and prompt/schema design as the leverage points, not model
breadth.
